% Modernised reading — fonds Alexandre Grothendieck, shelfmark 29, batch 11
% (pages 201-216). Last batch of the shelfmark.
%
% This is an INTERPRETATION, not a transcription. It re-reads the pages in
% today's notation and vocabulary, states the hypotheses the manuscript leaves
% implicit, and drops the critical apparatus so that the mathematics reads
% straight through. Everything the brackets and underlines carried moves into a
% footnote. In any dispute of substance the transcription governs: whoever
% cites Grothendieck must cite batch-11.fr.tex.
%
% Scope: the folder was read WHOLE before any of this was written — all eleven
% batches, pages 1-216, in one pass — and one file was then written per batch.
% This file is where the folder closes on itself, and the closure is only
% visible at folder scope: the pair (C, s) of pages 201-203, presented here as
% a « famille génératrice » of the champ R, IS the donnée de ramification of
% 1967; the cartesian sections of page 202 ARE the revêtements de type R of
% page 13; the formula of page 203 IS the colimit-of-limits of page 13; and
% the frames construction of page 209 IS Serre's proof of page 152. Each
% junction is named where it occurs.
%
% Departures from the page, each footnoted where it occurs:
%   - the equivalence of page 203 is CONDITIONAL: full faithfulness of alpha is
%     proved on pages 203-206, essential surjectivity is announced on page 206
%     and the proof does not follow. The reading says so and does not assert
%     the equivalence;
%   - page 213's dimension bound is written >= n - card(E)d where its own
%     equality case requires n - (card E - 1)d;
%   - page 213 labels two of its three partition lines « en 3 » where one must
%     be « en 1 ». Left as the page has it, and flagged.
%
% Pass: Opus 5 (claude-opus-5), 2026-08-29 — first pass, unchecked against the
% pages by a human.
\documentclass[11pt,a4paper]{article}
% Le préambule commun à toutes les transcriptions du fonds Grothendieck.
%
% Il tient en une idée : **l'incertitude se compose, elle ne se résout pas.**
% Une transcription qui lisse un mot douteux a détruit la seule chose pour
% laquelle on la faisait — un lecteur doit pouvoir savoir, à chaque endroit,
% ce qui était sur le feuillet et ce qui a été ajouté. Les six macros qui
% suivent sont donc l'appareil critique, et non de la décoration.
%
% Les mêmes commandes sont reconnues par `scripts/render.mjs`, qui produit la
% vue de lecture du site. Une macro ajoutée ici doit l'être aussi là-bas,
% faute de quoi le rendu échouera bruyamment — ce qui est voulu : un rendu
% silencieusement faux est pire qu'un rendu absent.

\ProvidesPackage{grothendieck}[2026/08/08 Transcriptions du fonds Grothendieck]

\usepackage{amsmath,amssymb,amsthm}
\usepackage{mathrsfs}
\usepackage{stmaryrd}
% Les diagrammes commutatifs sont partout dans ce fonds : c'est la langue
% même de la géométrie algébrique de Grothendieck, et une transcription qui
% les rendrait en prose ne transcrirait rien.
\usepackage{tikz-cd}
% Français par défaut — c'est la langue des feuillets — mais l'anglais est
% chargé aussi : la traduction et le résumé sont des documents anglais, et la
% typographie française (espaces avant les deux-points) y serait une faute.
% Ces documents basculent par \selectlanguage{english} en tête de corps.
\usepackage[english,french]{babel}
\usepackage{fontspec}
\usepackage{geometry}
\geometry{a4paper,margin=25mm,marginparwidth=28mm}
\usepackage{marginnote}
\usepackage{xcolor}
% ulem, not soul. `soul`'s \st and \ul are built on a character-by-character
% scan that XeLaTeX's font machinery breaks: the document fails with an
% undefined control sequence rather than degrading. ulem does the same two jobs
% and is Unicode-engine safe. `normalem` stops it hijacking \emph, which the
% transcriptions use for Grothendieck's own emphasis.
\usepackage[normalem]{ulem}
\usepackage{hyperref}
\hypersetup{colorlinks=true,linkcolor=black,urlcolor=black,pdfborder={0 0 0}}

% --- Métadonnées du lot -----------------------------------------------------
% Reprises telles quelles par le script de rendu, qui les affiche en tête de
% la vue de lecture. Elles fixent ce que le fichier prétend couvrir : sans
% elles, une transcription flotte sans point d'attache dans un fonds de seize
% mille feuillets.

\newcommand{\folder}[1]{\gdef\grothFolder{#1}}
\newcommand{\batch}[1]{\gdef\grothBatch{#1}}
\newcommand{\leaves}[2]{\gdef\grothFirst{#1}\gdef\grothLast{#2}}
\newcommand{\foldertitle}[1]{\gdef\grothFoldertitle{#1}}
\gdef\grothFolder{?}\gdef\grothBatch{?}\gdef\grothFirst{?}\gdef\grothLast{?}\gdef\grothFoldertitle{}

% --- Le résumé --------------------------------------------------------------
%
% Il ouvre la lecture modernisée, sous le titre « Résumé », et il est composé
% en petit. Ce n'est pas un ornement : c'est le seul passage du document qui
% ne soit pas des mathématiques, et un lecteur doit pouvoir le reconnaître
% avant de l'avoir lu — puis le sauter s'il connaît déjà le sujet, ou s'y
% arrêter s'il ne le connaît pas. Le filet qui le suit marque la reprise du
% corps.

\newenvironment{resume}
  {\small\setlength{\parskip}{0.45em}}
  {\par\medskip\hrule\medskip}

% --- Mots-clés ---------------------------------------------------------------
%
% En anglais, en clôture du résumé de la lecture modernisée : le vocabulaire
% moderne sous lequel chercher ce que les feuillets construisent. Le manifeste
% du site (`npm run manifest`) les extrait pour étiqueter la cote entière —
% cette ligne est la source unique des tags, il n'y a pas de fichier à côté.
\newcommand{\keywords}[1]{\par\smallskip\noindent
  {\footnotesize\textsc{Keywords} --- \textit{#1}}\par}

% --- L'appareil critique ----------------------------------------------------

% Le numéro de feuillet, en marge. C'est l'ancre qui permet de régler un
% désaccord sur une formule en regardant *un* feuillet plutôt qu'une chemise
% de six cents. Le site s'en sert aussi pour tourner les pages du fac-similé
% quand on fait défiler la transcription.
\newcommand{\leaf}[1]{%
  \marginnote{\footnotesize\color{gray!70}\textsf{#1}}%
  \label{leaf:#1}\ignorespaces}

% Illisible : on ne devine pas. Un mot inventé qui se lit comme les autres est
% le pire résultat possible.
\newcommand{\ill}{\textcolor{red!60!black}{[\,\ldots\,]}}

% Lecture douteuse : le mot est donné, et signalé comme incertain.
\newcommand{\uncertain}[1]{\uline{#1}}

% Ajout de l'éditeur — une lettre manquante, un mot sous-entendu.
\newcommand{\add}[1]{\textcolor{blue!50!black}{[#1]}}

% Biffé par Grothendieck lui-même. Ce qu'il a rayé fait partie du document :
% une rature dit souvent où la pensée a changé de direction.
\newcommand{\struck}[1]{\sout{#1}}

% Note du transcripteur, sur l'état matériel du feuillet.
\newcommand{\note}[1]{\par\smallskip{\small\color{gray!60!black}\textit{[#1]}}\par\smallskip}

% Note marginale de Grothendieck, distincte du corps du texte.
\newcommand{\marginal}[1]{\par\smallskip\noindent\hspace{1em}%
  {\small\color{gray!60!black}#1}\par\smallskip}

% --- Titre ------------------------------------------------------------------

\AtBeginDocument{%
  \begin{center}
    {\large\bfseries Fonds Alexandre Grothendieck --- cote n\textsuperscript{o}~\grothFolder}\\[2pt]
    {\itshape\grothFoldertitle}\\[4pt]
    {\small feuillets \grothFirst--\grothLast\ (lot \grothBatch)}\\[2pt]
    {\footnotesize\color{gray!60!black}Transcription établie d'après le fac-similé numérisé
     par l'Université de Montpellier.\\
     Les lectures incertaines sont soulignées, les illisibles notées [\,\ldots\,],
     les ajouts entre crochets.}
  \end{center}
  \vspace{1em}}

\endinput


\folder{29}
\batch{11}
\pages{201}{216}
\dating{1959-1969}
\watermark{Édition de démonstration\\interprétation personnelle de l'œuvre}
\foldertitle{Groupe fondamental [Autour de SGA 4 et SGA 7] — lecture modernisée}

\begin{document}

\section*{Résumé}

\begin{resume}
Ces seize pages ferment le carton, et elles le ferment sur lui-même.

La première moitié reprend la question laissée ouverte au lot précédent. On a
maintenant un \emph{domaine de ramification} : un champ au-dessus du site
étale, avec sa réalisation géométrique. Comment le décrire ? En le
\emph{présentant} : on se donne une famille d'objets qui l'engendre, et l'on
demande si le champ se reconstitue à partir d'elle. Grothendieck écarte d'abord
la présentation naïve --- par des schémas à opérateurs $(Z_{i}, G_{i})$ --- avec
un jugement qui vaut d'être retenu : « la formulation de telles données est
pénible et finalement peu utile ». Il lui préfère une catégorie de modèles $C$
au-dessus du site étale et un foncteur $\varphi : C \to R$, et il pose
$s = r\varphi$.

Il faut s'arrêter là un instant, parce que le carton vient de boucler. Le
couple $(C, s)$ est \emph{exactement} la « donnée de ramification » que le
tapuscrit de 1967 posait à sa page 12, deux cents pages plus tôt. Les
« sections cartésiennes » que la page 202 construit sont \emph{exactement} les
revêtements de type $\mathcal{R}$ de la page 13. Et la formule que la page 203
vise --- une limite inductive sur les cribles de limites projectives sur chaque
crible --- est \emph{exactement} celle qui définissait
$\mathrm{Rev}^{\mathcal{R}}(S)$. Ce qui était en 1967 une \emph{définition}
posée est devenu ici un \emph{théorème de présentation} : la donnée de 1967
n'était pas l'objet, c'était une présentation de l'objet, et le carton l'aura
mis dix ans à le démontrer.

Le mot \emph{démontrer} est trop fort, et il faut le dire : la démonstration
est faite à moitié. Que le foncteur soit pleinement fidèle occupe quatre pages,
et l'argument est joli --- il repose sur le fait qu'un objet subordonné à un
modèle est recouvert par les images de ses sections, ce qui permet de remplacer
un objet quelconque par le modèle lui-même. Puis la page 206 écrit « Prouvons
que $\alpha$ est essentiellement surjectif », et la démonstration ne suit pas.
La suite est d'un autre sujet.

La seconde moitié est le plus lisible de tout le carton, écrit au net à l'encre
bleue, et c'est l'axiomatique des \emph{catégories multigaloisiennes} --- ce
mot que la page 30 employait déjà en 1967 comme s'il allait de soi. Une
catégorie galoisienne, c'est celle des ensembles finis où opère un groupe
profini. Une catégorie multigaloisienne, c'est la même chose avec
\emph{plusieurs} composantes : un produit de catégories galoisiennes, une par
composante connexe de l'objet final --- donc un \emph{groupoïde} et non un
groupe. Trois caractérisations sont données, et la première est la plus
frappante : ce sont les objets localement constants finis d'un topos. Les six
axiomes G0 à G5 en donnent la version élémentaire, et Grothendieck note du
dernier qu'il le soupçonne redondant, sans le démontrer.

Deux feuillets insérés font autre chose et méritent d'être signalés : ils
axiomatisent directement la notion de \emph{morphisme de degré $n$} --- ce que
veut dire, sans jamais compter de points, qu'un revêtement a $n$ feuillets ---
et démontrent qu'une telle notion, si elle existe, est unique.

Le carton s'arrête sur une proposition sans démonstration, à la moitié d'une
page.

\keywords{multigaloisian category, Galois topos, fundamental groupoid, generating family, degree of a morphism}
\end{resume}

\pagerange{201}{202}
\section*{Familles génératrices, et la donnée de 1967 retrouvée (pages 201 et 202)}

Grothendieck pagine cette suite 11 à 16, à la suite du 1--10 du lot précédent.

\subsection*{Ce qu'est une famille génératrice}

Soit $(R, r)$ un domaine de ramification sur $S$. Une famille d'objets
$X_{i}/S_{i}$ de $R$, avec $S_{i} \in \mathrm{Et}(S)$, est dite
\emph{génératrice} si, pour tout objet $X$ de $R$ sur un $S'$, les $S''$ étales
sur $S'$ pour lesquels il existe un $i$, un morphisme $S'' \to S_{i}$ et une
subordination de $X_{S''}$ à $X_{i,S''}$, recouvrent $S'$.

En clair : \emph{tout revêtement est localement subordonné à un des
générateurs}. C'est la même condition que le crible couvrant de 1967, dite du
côté du champ.

\subsection*{Le refus de la présentation naïve}

On pourrait vouloir reconstituer $(R, r)$ à partir des seuls
$(Z_{i}, G_{i}) = (r(X_{i}), \operatorname{Aut}(X_{i}))$, en précisant les
compatibilités. Grothendieck écarte ce parti, et son jugement mérite d'être
retenu tel quel : « la formulation de telles données est \emph{pénible et
finalement peu utile} ».\footnote{La phrase de la page se défait entre « Si de
plus » et « les $Z_{i}$, $G_{i}$ doivent permettre » ; ce qui se lit est le
projet et son abandon. C'est le troisième refus de ce genre dans le carton,
après le « voile pudique » de la page 27 et le « pénible » implicite des cinq
conditions a) à e) de la page 32 --- et les trois portent sur le même objet :
l'explicitation en coordonnées de ce qui se dit sans coordonnées.}

\subsection*{Le couple $(C, s)$, et où le carton se referme}

Il lui préfère d'associer à $R$ une catégorie $C$ sur $\mathrm{Et}(S)$ et un
foncteur
\[
  \varphi : C \longrightarrow R ,
\]
dit \emph{générateur} si les $\varphi(M)$ forment une famille génératrice ---
c'est-à-dire si, pour tout $X \in R$ sur $S'$, les $S''$ étales sur $S'$ tels
que $X_{S''}$ soit subordonné à un objet de $C_{S''}$ recouvrent $S'$. Et il
pose
\[
  s \;=\; r\varphi \;:\; C \longrightarrow \mathrm{Fl}_{S} ,
\]
en annonçant qu'on peut alors, à équivalence près, reconstituer $(R, r)$ en
termes de $(C, s)$ --- et de même les $(R_{S'}, r_{S'})$, avec toutes les
fonctorialités.

\begin{quote}
\textbf{Le couple $(C, s)$ est la donnée de ramification de 1967.} Le tapuscrit
de la page 12 posait : « une donnée de ramification sur $S$ est un couple
$(\mathcal{C}, s)$, où $\mathcal{C}$ est une catégorie fibrée sur
$\mathrm{Et}(S)$ et $s : \mathcal{C} \to \mathrm{Fl}_{S}$ un foncteur
cartésien ». C'est le même couple, avec le même nom pour le second membre. Ce
qui a changé est son statut : en 1967 c'était la donnée dont on \emph{déduisait}
la catégorie ; ici c'est une \emph{présentation} de la catégorie, qui existe
d'abord, et dont on demande si elle la détermine.
\end{quote}

\subsection*{Les sections cartésiennes}

Soit $X$ un objet de $R_{S}$, et $C_{X}$ la sous-catégorie pleine de $C$ formée
des objets $M$ sur un $S'$ tels que $X_{S'}$ soit subordonné à $\varphi(M)$ et
que $\varphi(M) \to X_{S'}$ soit un épimorphisme. Pour un tel $M$,
\[
  r\bigl(X_{S'} \times \varphi(M)\bigr) \longrightarrow r\varphi(M) = s(M)
\]
est un \emph{revêtement étale}. Donc
\[
  M \longmapsto r\bigl(X \times_{S'} \varphi(M)\bigr) \big/ s(M)
\]
donne une \emph{section cartésienne} de la catégorie fibrée des revêtements
étales relatifs, tirée en arrière sur $C_{X}$ par
$C_{X} \xrightarrow{\ s\ } \mathrm{Fl}_{S}$.

\begin{quote}
\textbf{Ce sont les revêtements de type $\mathcal{R}$ de 1967.} La page 13 du
carton posait : « on appellera revêtement de type $\mathcal{R}$ de $S$ une
section cartésienne de $\mathcal{D}$ sur $\mathcal{C}$ », $\mathcal{D}$ étant
l'image inverse par $s$ de la catégorie fibrée des revêtements étales finis.
C'est mot pour mot la même construction, dans l'autre sens : là on partait de
$\mathcal{C}$ pour \emph{fabriquer} les sections cartésiennes, ici on part d'un
objet $X$ du champ pour lui \emph{associer} une section cartésienne.
\end{quote}

\subsection*{Les sous-catégories $R_{S}^{C_{0}}$}

Pour un crible \emph{admissible} $C_{0}$ de $C$, soit $R_{S}^{C_{0}}$ la
sous-catégorie \emph{pleine} de $R_{S}$ formée des objets $X$ tels que, pour
tout $S' \in \mathrm{Et}(S)$, $X_{S'}$ soit localement subordonné à \emph{tout}
objet de $C_{0}$ sur $S'$. C'est une sous-catégorie multigaloisienne de
$R_{S}$, stable par limites projectives et par sommes directes, et l'on
dispose d'un foncteur
\[
  (\ast) \qquad
  \alpha : R_{S}^{C_{0}} \longrightarrow
  \varprojlim_{C_{0}} \bigl(\mathcal{E}\,|\,C_{0}\bigr) ,
\]
$\mathcal{E}$ désignant la catégorie fibrée des revêtements étales
relatifs.\footnote{Le symbole que la page emploie pour cette catégorie fibrée
est un E de fantaisie qu'elle ne nomme nulle part ; la transcription le rend
$\mathcal{E}$ partout où il paraît, et cette lecture fait de même. C'est le
$\mathcal{D}$ de la page 12, pris au-dessus de $C_{0}$.}

\pagerange{203}{206}
\section*{$\alpha$ est pleinement fidèle ; et la suite manque (pages 203 à 206)}

\subsection*{La formule visée}

Faisons parcourir à $C_{0}$ les cribles admissibles de $C$, et supposons la
catégorie fibrée $C$ \emph{localement filtrante} --- ce qui entraîne que
l'intersection de deux cribles couvrants est couvrante. Alors les
$R_{S}^{C_{0}}$ forment une famille \emph{filtrante} de sous-catégories
multigaloisiennes pleines de $R_{S}$, de réunion $R_{S}$. Donc, si les
$\alpha_{0}$ sont des \emph{équivalences}, on aura
\[
  R_{S} \;\simeq\; \varinjlim_{C_{0}} \; \varprojlim_{C_{0}}
  \bigl(\mathcal{E}\,|\,C_{0}\bigr) ,
\]
formule qui explicite $R_{S}$ en termes de $(C, s)$.

\begin{quote}
\textbf{C'est la formule de 1967.} La page 13 du carton posait
$\mathrm{Rev}^{\mathcal{R}}(S) = \varinjlim_{\mathcal{C}_{0}}
\mathrm{Rev}^{\mathcal{R}}_{\mathcal{C}_{0}}(S)$, où
$\mathrm{Rev}^{\mathcal{R}}_{\mathcal{C}_{0}}(S)$ était la catégorie des
sections cartésiennes au-dessus de $\mathcal{C}_{0}$, c'est-à-dire précisément
$\varprojlim_{\mathcal{C}_{0}}$. Une limite inductive sur les cribles, de
limites projectives sur chaque crible : c'est la même formule, aux deux
extrémités du carton. Et les deux hypothèses (H1) et (H2) que la lecture du lot
1 avait dû isoler --- la pleine fidélité des foncteurs de transition,
l'existence de $r$ --- sont maintenant \emph{ce qu'il s'agit de démontrer}.
\end{quote}

Et l'on annonce que pour cela il suffit que $\varphi$ soit \emph{pleinement
fidèle}.

\subsection*{Ce qui est facile}

$\alpha$ est un foncteur exact de catégories multigaloisiennes ; il suffit donc,
pour la fidélité, de voir que
\[
  \alpha(X) = \varnothing \;\Longrightarrow\; X = \varnothing .
\]
C'est immédiat : la question est locale sur le site étale, il existe un
$M \in C_{0,S}$, et
$r(\varphi(M) \times X) = \varnothing$ entraîne
$\varphi(M) \times X = \varnothing$ par la proposition 2 de la page 192.

\subsection*{Ce qui demande du travail : la pleine fidélité}

Soit $X \in R_{S}^{C_{0}}$ et $Z = \alpha(X) = Z_{1} \sqcup Z_{2}$ ; il faut
montrer que cette décomposition provient d'une décomposition de $X$.

La question est locale sur $S$ ; on peut donc supposer qu'il existe
$M \in C_{S}$ avec $\varphi(M) \to e$ épimorphique. Considérons
\[
  r\bigl(X \times \varphi(M)\bigr) \;=\; \alpha(X)(M) \;=\; Z(M)
  \longrightarrow r(\varphi(M)) = s(M) .
\]
Il se coupe en $Z_{1}(M) \sqcup Z_{2}(M)$, donc \emph{par l'axiome R2}
$X \times \varphi(M)$ se coupe en $X_{1} \sqcup X_{2}$.

Reste à montrer que cette décomposition provient d'une décomposition de $X$ :
c'est-à-dire que pour tout objet $T$ de $R_{S}^{C_{0}}$ et deux morphismes
$T \rightrightarrows \varphi(M)$, les deux images inverses de $X_{1}$
coïncident.

\begin{tikzcd}[column sep=small, row sep=small]
T \arrow[d, bend left=12] \arrow[d, bend right=12]
  & X \times T \arrow[d, bend left=12] \arrow[d, bend right=12] & \\
\varphi(M) \arrow[d]
  & X \times \varphi(M) \arrow[l, no head] \arrow[r] \arrow[d] & X_1 \\
e & X \arrow[l, no head] &
\end{tikzcd}

L'argument est le suivant, et c'est le morceau de bravoure du lot.

Le morphisme $T \times \varphi(M) \to \varphi(M)$ est un épimorphisme. Or
$T \times \varphi(M)$, étant subordonné à $\varphi(M)$, est \emph{constant par
morceaux} sur $\varphi(M)$ ; et un objet constant par morceaux et épimorphique
est \emph{recouvert}, pour la topologie canonique, par les images des
\emph{sections} $w$ de $T \times \varphi(M)$ sur $\varphi(M)$. On peut donc
remplacer le couple $T \rightrightarrows \varphi(M)$ d'abord par les composés
\[
  T \times \varphi(M) \xrightarrow{\ \mathrm{pr}_{1}\ } T
  \rightrightarrows \varphi(M)
\]
--- ce qui est licite, $\varphi(M) \to e$ étant un épimorphisme --- puis par
les composés $\varphi(M) \xrightarrow{\ w\ } T \times \varphi(M) \to \cdots$.
\emph{On est ainsi ramené au cas $T = \varphi(M)$.} Mais alors, $\varphi$ étant
pleinement fidèle, les deux flèches $M \rightrightarrows M$ sont identiques.

Le mécanisme mérite d'être nommé : \emph{quand tout objet est localement
constant sur un modèle, on peut remplacer un objet quelconque par le modèle
lui-même, parce que les sections le recouvrent}. C'est exactement pour cela que
la subordination avait été définie.

On obtient donc une décomposition
\[
  X \;=\; X_{1} \sqcup X_{2}
\]
uniquement déterminée par la condition $Z_{1} = \alpha(X_{1})$. Il reste à voir
qu'elle ne dépend pas du choix de $M$ : c'est ce que fait la fin de la page 205
et la page 206, en localisant sur $S$ et en utilisant que $C$ est localement
filtrante.\footnote{Cette dernière partie est, sur la page 205, accolée à
gauche et barrée de deux longues diagonales ; elle se lit par endroits et par
endroits non. Le principe --- comparer les décompositions définies par $M$ et
par $M'$ en passant à un $M''$ qui les majore, ce que la filtration locale
fournit --- est ce que les phrases portantes autorisent ; le détail ne l'est
pas, et il n'est pas reconstruit ici.}

\subsection*{Et la suite manque}

La page 206 se termine sur une seule ligne :

\begin{quote}
Prouvons que $\alpha$ est \emph{essentiellement surjectif}.
\end{quote}

\emph{La démonstration ne suit pas.} La page s'arrête là et la suite du carton
est d'un autre sujet.

Il faut donc dire ce qui est acquis et ce qui ne l'est pas. Est acquise la
pleine fidélité de $\alpha$ --- $R_{S}^{C_{0}}$ s'envoie fidèlement et
pleinement dans la limite projective. N'est pas acquise l'essentielle
surjectivité, donc \emph{l'équivalence annoncée à la page 203 reste
conditionnelle}, et avec elle la reconstitution de $(R, r)$ à partir de
$(C, s)$. La boucle que le carton ferme --- de la donnée de 1967 au champ, et
du champ à la donnée --- est fermée dans un sens et pas dans l'autre.

\pagerange{208}{209}
\section*{Catégories multigaloisiennes : les conditions A et B (pages 208 et 209)}

À partir d'ici et jusqu'à la fin, le texte est écrit au net, à l'encre bleue,
et se lit presque de bout en bout --- à la différence de tout ce qui précède
dans ce carton. Grothendieck pagine 1 à 4, insère deux feuillets « 1 bis » et
« 1 ter », puis 5 et 6.

\subsection*{A) Ce que c'est}

Soit $\mathcal{U}$ un univers et $C$ une catégorie. Les conditions suivantes
sont équivalentes.

\begin{enumerate}
  \item[(i)] Il existe un $\mathcal{U}$-topos $\mathcal{T}$ et une équivalence
    $C \simeq \mathcal{T}_{\mathrm{cf}}$, où $\mathcal{T}_{\mathrm{cf}}$ est la
    sous-catégorie pleine des objets \emph{localement constants
    finis}.\footnote{L'indice est « cf », et c'est Grothendieck lui-même qui le
    glose à la ligne au-dessus : « constants finis ». La transcription le
    signale ; c'est cette glose, et non le tracé, qui fixe la lecture.}
  \item[(ii)] $C$ est une $\mathcal{U}$-catégorie dont l'ensemble des classes
    d'isomorphie d'objets est $\mathcal{U}$-petit, et il existe sur $C$ une
    topologie --- moins fine que la topologie canonique --- telle que
    $C \to \tilde{C}$ induise une équivalence
    $C \xrightarrow{\ \sim\ } \tilde{C}_{\mathrm{cf}}$.
  \item[(ii bis)] La même chose, en prenant la topologie \emph{la plus fine de
    la descente effective universelle}.
  \item[(iii)] Pour tout objet $X$ de $C$, les objets $S$ tels que $X_{S}$ soit
    \emph{constant fini} sur $S$ forment une famille de descente effective
    universelle.
  \item[(iv)] Quatre conditions, portées sur un feuillet inséré (page 212).
\end{enumerate}

La condition (i) est la définition conceptuelle : \emph{une catégorie
multigaloisienne est la catégorie des systèmes locaux finis d'un topos}. La
condition (iii) est la définition opératoire : \emph{tout objet devient
constant après un changement de base convenable}. Que les deux coïncident est
tout le contenu de A.

\subsection*{B) Le cas multigaloisien proprement dit}

Les conditions suivantes sont équivalentes.

\begin{enumerate}
  \item[(i)] $C$ satisfait aux conditions de A), et son objet final est somme
    \emph{effective universelle} d'objets connexes.
  \item[(ii)] $C$ satisfait aux axiomes \textbf{G0} à \textbf{G5} ci-dessous.
  \item[(iii)] $C$ est équivalente à une catégorie produit
    $\prod_{i \in I} C_{i}$, les $C_{i}$ étant \emph{galoisiennes}.
\end{enumerate}

Avec la note qui explique le mot : \emph{les $C_{i}$ correspondent aux
composantes connexes $e_{i}$ de l'objet final} ; et la donnée d'une équivalence
$C_{i} \simeq C(\pi_{i})$ avec la catégorie des ensembles finis à opérateurs
d'un groupe profini $\pi_{i}$ correspond à la donnée d'un \emph{foncteur fibre}
de $C_{/e_{i}}$, c'est-à-dire d'un « revêtement universel » de $e_{i}$.

\begin{quote}
Voilà ce que « multigaloisien » veut dire, et c'est la réponse à une question
que le carton posait à sa page 30, en 1967, en employant le mot sans le
définir. Une catégorie galoisienne est celle des ensembles finis sur lesquels
opère \emph{un} groupe profini ; elle exige que l'objet final soit connexe. Une
catégorie multigaloisienne est un \emph{produit} de catégories galoisiennes,
une par composante connexe de l'objet final : elle correspond donc non à un
groupe mais à un \emph{groupoïde} profini, et c'est ce qu'il faut quand on ne
veut pas supposer la base connexe. C'est exactement ce qui manquait au
paragraphe 2 du tapuscrit de 1967, dont le groupe
$\pi_{1}^{\mathcal{R}}(S,s)$ n'existait que sous une hypothèse de connexité du
site.
\end{quote}

\subsection*{Les axiomes G0 à G5}

\begin{itemize}
\item[\textbf{G0}] $C$ est une $\mathcal{U}$-catégorie, et l'ensemble des
  classes d'isomorphie d'objets est $\mathcal{U}$-petit.
\item[\textbf{G1}] Existence des limites projectives \emph{finies} et des
  sommes \emph{finies}.
\item[\textbf{G2}] Les relations d'équivalence sont \emph{effectives}, et les
  épimorphismes sont universels.
\item[\textbf{G3}] Tout morphisme $X \to Y$ se factorise en un
  \emph{épimorphisme} suivi d'un \emph{isomorphisme sur un facteur direct}.
\item[\textbf{G4}] L'objet final est somme \emph{effective universelle}
  d'objets connexes.
\item[\textbf{G5}] Soit $f : X \to Y$ avec $X$ et $Y$ connexes. Alors il existe
  un entier $N$ tel que, pour tout $Y' \to Y$ non vide, \dots\footnote{L'énoncé
  de G5 est le seul de la liste que la page laisse incomplet : la conclusion
  est biffée et réécrite, et la transcription ne la lit pas. Ce qu'il vise se
  déduit de l'usage : c'est une borne uniforme sur le nombre de composantes
  connexes de $X \times_{Y} Y'$, c'est-à-dire une condition de \emph{degré
  borné} qui interdit à un morphisme entre objets connexes de se scinder
  indéfiniment. La page 211 revient dessus, et pour en dire qu'elle le soupçonne
  redondant.}
\end{itemize}

L'axiome \textbf{G3} est le seul de la liste qui ait traversé tout le carton :
c'est la propriété que la page 24, en 1967, obtenait comme un \emph{théorème}
sur $\mathrm{Rev}^{\mathcal{R}}(S)$ --- « tout morphisme de la catégorie se
factorise en le produit d'un épimorphisme et d'un monomorphisme \dots\ tout
monomorphisme est un isomorphisme avec un sommande direct » --- et c'est
l'axiome G3 de SGA 1, exposé V.

\pagerange{209}{211}
\section*{Les démonstrations, et un axiome soupçonné (pages 209 à 211)}

\subsection*{Démonstration de A : le revêtement des repères, une seconde fois}

Les implications (iii) $\Rightarrow$ (ii bis) $\Rightarrow$ (ii) $\Rightarrow$
(i) sont triviales. Pour (i) $\Rightarrow$ (iii) :

Soit $X \in \mathcal{T}_{\mathrm{cf}}$. L'objet final de $\mathcal{T}$ est somme
d'objets $e_{i}$ ($i \in \mathbb{N}$) au-dessus desquels $X|e_{i}$ est de rang
$i$. Posons alors
\[
  P_{i} \;=\; \mathbf{Isom}_{e_{i}}\bigl([1,i]_{e_{i}},\ X\bigr) ,
\]
le faisceau des identifications de $X$ avec l'ensemble constant à $i$ éléments.
C'est un objet de $\mathcal{T}_{\mathrm{cf}}$, l'application
$P_{i} \to e_{i}$ est un épimorphisme --- autrement dit $P_{i}$ est un
\emph{torseur sous $(\mathfrak{S}_{i})_{e_{i}}$} --- et
\[
  X_{P_{i}} \;\simeq\; [1, i]_{P_{i}} .
\]
Les $P_{i} \to e$ forment donc une famille de descente effective universelle
pour $\mathcal{T}$, donc aussi pour $\mathcal{T}_{\mathrm{cf}}$, qui est stable
par limites projectives et inductives finies.

\begin{quote}
\textbf{C'est la démonstration de Serre, page 152 du carton.} Là, $X \to Y$
était un morphisme fini étale de rang $n$, et $P$ le complémentaire de la
diagonale dans $X^{n}$ --- c'est-à-dire le schéma des $n$-uplets de points deux
à deux distincts, ou encore le schéma des identifications de la fibre avec
$I_{n}$ ---, torseur sous $\mathfrak{S}_{n}$. Ici, $\mathbf{Isom}([1,i], X)$
est le même objet, écrit dans un topos quelconque. Le carton contient donc la
construction dans son cas géométrique et dans sa forme abstraite, à
soixante pages de distance, et l'attribution à Serre est portée sur la
première.
\end{quote}

\subsection*{Démonstration de B, et l'axiome G5}

Les conditions A(ii) impliquent aisément G0, G1, G2, G3. Pour G4, il faut
prouver que pour tout $X \in C$ les $S_{i} \to e$ tels que $X_{S_{i}}$ soit
constant fini recouvrent $e$ pour la topologie de la descente effective ;
cela s'obtient en utilisant G2 et G4, en sommant sur les composantes connexes
$e_{i}$ de l'objet final.

Puis l'énoncé de récurrence : si $X \to Y$ est tel que $Y$ soit non vide et $X$
connexe, alors il existe $Y' \to Y$ de descente effective universelle avec
$X_{Y'} \simeq [1,n]_{Y'}$ ; on procède par récurrence sur le rang, et c'est
là que G5 sert.

\textbf{Remarque.} Grothendieck note qu'il soupçonne \textbf{G5} d'être
\emph{redondant}, c'est-à-dire conséquence des autres axiomes, et indique ce
qu'il faudrait pour le montrer : pour toute composante connexe $e_{i}$ de
l'objet final, l'existence d'un foncteur exact
\[
  F : C \longrightarrow (\text{Ens.\ finis}) \qquad \text{tel que}
  \qquad F(e_{i}) \neq \varnothing .
\]
Il ne le démontre pas.\footnote{C'est le second axiome que le carton
soupçonne de trop, après celui de la page 157 --- « $F(X) = \varnothing
\Rightarrow X = 0$ » --- qui, lui, était biffé parce que \emph{démontré}
redondant, conséquence de Gal 5. Ici le soupçon reste un soupçon. Une bonne
partie de la page 211 est barrée de longues diagonales et ne se lit que par
fragments ; ce qui précède est ce que les phrases portantes autorisent.}

\pagerange{212}{214}
\section*{Compter les feuillets sans compter de points : les morphismes de degré $n$ (pages 212 et 214)}

Deux feuillets insérés, que Grothendieck numérote « 1 bis » et « 1 ter » dans
un cercle, portent les quatre conditions que la page 208 annonçait sous A(iv).
Elles sont numérotées à rebours --- (3), (2), (1), (0) --- chaque chiffre
cerclé.

Ce qu'elles font mérite d'être dit avant d'être lu : elles axiomatisent
directement la notion de \emph{morphisme de degré $n$}. Dire qu'un revêtement a
$n$ feuillets suppose d'ordinaire qu'on compte les points d'une fibre ; ici, il
n'y a pas de points, et il faut donc dire ce qu'un degré est par ses
propriétés seules.

\begin{itemize}
\item[\textbf{(3)}] Pour tout objet non vide, on se donne, pour chaque $n$, une
  classe de « morphismes de degré $n$ », telle que :
  \begin{enumerate}
    \item[a)] elle soit stable par changement de base ;
    \item[b)] elle soit \emph{additive} pour les sommes finies ;
    \item[c)] un isomorphisme soit de degré 1 ;
    \item[d')] pour tout $f : X \to Y$, $Y$ soit somme directe disjointe
      effective universelle $Y = \coprod_{n \in \mathbb{N}} Y_{n}$ avec
      $X_{n} = X \times_{Y} Y_{n} \to Y_{n}$ de degré $n$ ;
    \item[d)] si $f$ est de degré $n \geqslant 1$, $f$ soit un morphisme de
      descente \emph{effective} --- donc universelle, grâce à a) ;
    \item[e)] si $f$ est de degré $0$, alors $X = \varnothing$, et $f$ soit un
      isomorphisme sur un \emph{sommande direct universel}.
  \end{enumerate}
\item[\textbf{(2)}] Tous les monomorphismes sont d'un type prescrit.
\item[\textbf{(1)}] Existence des limites projectives \emph{finies}.
\item[\textbf{(0)}] $C$ est une $\mathcal{U}$-catégorie et ses classes
  d'isomorphie forment un ensemble $\mathcal{U}$-petit.
\end{itemize}

La condition d') est celle qui fait tout le travail : elle dit que le degré
n'est pas un nombre attaché au morphisme, mais une \emph{décomposition de la
base}. C'est la traduction exacte, sans points, du fait qu'un morphisme fini
étale a un rang localement constant.

\subsection*{La récurrence}

De ces conditions on déduit A(iii) : pour $X$ donné, les $S \to e$ tels que
$X_{S}$ soit constant fini forment une famille de descente effective
universelle. On se ramène, par a) et d'), au cas d'un $X \to Y$ de degré $n$.
Si $n = 0$ on conclut par 3 e), $X$ étant vide et $Y \to Y$ convenant. Si
$n > 0$, on procède par récurrence, l'énoncé étant supposé acquis pour les
entiers $< n$ :

Utilisant 3 d), posons $X' = X \times_{Y} X$ au-dessus de $Y' = X$ :

\begin{tikzcd}[column sep=small, row sep=small]
X \arrow[d] & X \times_Y X = X' \arrow[l] \arrow[d] \\
Y & X = Y' \arrow[l, no head]
\end{tikzcd}

La \emph{diagonale} $\Delta_{X,Y}$ existe par (1), et par (2) c'est un
isomorphisme sur un sommande direct universel : donc
$X' \simeq Y' \sqcup X'_{1}$. Par 3 a), $X' \to Y'$ est de degré $n$ ; par 3 c),
$Y' \to Y'$ est de degré 1 ; donc par additivité, $X'_{1} \to Y'$ est de degré
$n - 1$. L'hypothèse de récurrence fournit alors une famille
$Y''_{s} \to Y'$ de descente effective universelle qui trivialise $X'_{1}$ ;
et comme $X' \simeq Y' \sqcup X'_{1}$ \emph{universellement},
\[
  X'' \;=\; X \times_{Y} Y'' \;\simeq\; Y'' \sqcup
  \underbrace{\bigl(X'_{1} \times_{Y'} Y''\bigr)}_{\textstyle
  Y'' \times [1, n-1]} \;=\; Y'' \times [1, n] .
\]
Et $Y'' \to Y$ est de descente effective universelle comme composé de deux
tels.

\subsection*{L'unicité}

\textbf{Conclusion.} Dans les conditions de A(iv), la notion de « morphisme de
degré $n$ » est \emph{uniquement déterminée}, et elle équivaut à : il existe
$Y' \to Y$ épimorphisme effectif universel, et de descente effective, avec
\[
  X \times_{Y} Y' \;\simeq\; [1, n]_{Y'} .
\]

C'est-à-dire : \emph{un morphisme est de degré $n$ s'il devient l'ensemble
constant à $n$ éléments après un changement de base convenable}, et cette
définition-là est forcée. On ne pouvait pas la prendre pour définition dès le
départ --- il aurait fallu savoir que « $n$ éléments » a un sens --- et c'est
pour cela que les six propriétés a) à e) sont posées d'abord.

Une \emph{explicitation} accompagne : pour vérifier qu'un $f$ est un morphisme
de descente effective, il suffit de savoir que c'est un épimorphisme effectif
universel, et que les relations d'équivalence sont effectives universelles.

\begin{tikzcd}[column sep=small, row sep=small]
Z'' \arrow[d]
  & Z \arrow[l, dashed] \arrow[d]
  & D \arrow[l, bend left=12] \arrow[l, bend right=12] \arrow[d] \\
Y & X \arrow[l]
  & X \times_Y X \arrow[l, bend left=12] \arrow[l, bend right=12]
\end{tikzcd}

\pagerange{213}{213}
\section*{Un feuillet d'un autre sujet : puissances fibrées et partitions (page 213)}

Une feuille en travers, sans numéro de sa main, barrée de deux longues
diagonales, et sans continuité avec ce qui l'entoure. Elle est transcrite à son
rang et lue ici de même.

Soit $f : X \to Y$ et $E$ un ensemble fini ; soit $X^{E}$ la puissance fibrée
de $X$ au-dessus de $Y$ indexée par $E$. À toute partition $P$ de $E$ est
associée une \emph{immersion fermée}

\begin{tikzcd}[column sep=large]
X^P \arrow[r, hook, "\alpha_P"] & X^E
\end{tikzcd}

--- celle qui identifie les coordonnées à l'intérieur de chaque bloc --- et
l'on a la décomposition disjointe
\[
  X^{E} \;=\; \bigcup_{P} \alpha_{P}\bigl((X^{P})^{\ast}\bigr) ,
\]
où $(X^{P})^{\ast}$ est la partie ouverte des points à coordonnées deux à deux
distinctes. De même, pour la petite diagonale de $Y$,
\[
  (f^{E})^{-1}\bigl(\Delta_{y}^{(E)}\bigr) \;=\;
  \bigcup_{P} \alpha_{P}(F_{P}) \qquad \text{(réunion disjointe)} ,
\]
les morceaux étant de dimension
\[
  n - (\operatorname{card} P - 1)\, d \;\geqslant\;
  n - (\operatorname{card} E - 1)\, d ,
\]
avec égalité pour la seule partition maximale
$P_{0}$.\footnote{La page écrit la minoration sous la forme
$\geqslant n - \operatorname{card}(E)\,d$, ce qui est vrai mais non optimal et
incompatible avec le cas d'égalité qu'elle énonce dans la même phrase : pour
$P = P_{0}$, la partition en $\operatorname{card} E$ blocs, la formule donne
$n - (\operatorname{card} E - 1)d$. C'est cette borne qui est écrite ici. Les
lettres $n$ et $d$ ne sont pas définies sur la page ; le cas travaillé
ci-dessous les fixe : $n = \dim X$ et $d$ la dimension relative de $f$.}

\subsection*{Le cas $\operatorname{card} E = 3$}

Les partitions de $E$ à trois éléments sont : celle en un bloc (la petite
diagonale), les trois en deux blocs, et celle en trois blocs. D'où
\[
  (f^{3})^{-1}\bigl(\Delta_{y}^{(3)}\bigr) \;=\;
  \underset{n - 2d}{F_{3}} \;\cup\;
  \underset{n - d}{\bigcup_{1 \leqslant i < j \leqslant 3} \alpha_{ij}(F_{2})}
  \;\cup\; \underset{n}{\Delta_{X}^{(3)}} ,
\]
les dimensions notées étant les dimensions « génériques ». Le cas numérique
travaillé sur la page est $\dim X = 3$, $\dim Y = 4$, avec
$\dim F_{3} = 1$ et $\dim \alpha_{ij}(F_{2}) = 2$ --- ce qui donne $n = 3$ et
$d = 1$.\footnote{La page numérote deux de ses trois lignes de partitions
« en 3 », là où l'une des deux doit porter « en 1 » : c'est la partition en un
seul bloc, qui donne la petite diagonale. Le chiffre est écrit deux fois de la
même main ; la transcription le conserve avec un \emph{sic}, et cette lecture
donne les trois lignes dans leur ordre correct.}

\textbf{N.B.} Si $f$ est convenable, $\overline{F}_{3} \cap
\Delta_{X}^{(3)} = \varnothing$ et
$\overline{\alpha_{ij}(F_{2})} \cap \Delta_{X}^{(3)} = \varnothing$. Mais la
question que la page pose et laisse ouverte est celle des intersections
croisées :
\[
  \overline{F}_{3} \cap \overline{\alpha_{ij}(F_{2})} ,
  \qquad
  \overline{F}_{3} \cap \alpha_{ij}(\overline{F}_{2}) \;?
\]

\pagerange{215}{216}
\section*{Le cas galoisien, et la dernière proposition (pages 215 et 216)}

\subsection*{C) Le cas galoisien}

Conditions équivalentes, définissant les catégories \emph{galoisiennes} :

\begin{enumerate}
  \item[(i)] Comme A(i), avec l'objet final de $\mathcal{T}$ \emph{connexe}.
  \item[] \textbf{(ii), (ii bis), (iii), (iv)} Comme dans A, en supposant de
    plus l'objet final de $C$ connexe.
  \item[(v)] Les axiomes \textbf{G0}, \textbf{G1}, \textbf{G2}, \textbf{G3},
    et de plus : il existe un foncteur
    \[
      F : C \longrightarrow (\text{Ens.\ fini})
    \]
    \emph{conservatif} et \emph{exact} --- respectant les limites projectives
    finies et commutant aux sommes finies.
  \item[(vi)] Il existe un \emph{groupe profini} $\pi$ et une équivalence
    $C \simeq C(\pi)$.
\end{enumerate}

C'est la liste de SGA 1, exposé V, dans sa forme définitive : les axiomes de la
catégorie, l'existence d'un foncteur fibre exact et conservatif, et la
conclusion --- la catégorie est celle des ensembles finis à opérateurs d'un
groupe profini. La différence avec B tient en un mot : ici l'objet final est
connexe, et le groupoïde se réduit à un groupe.

Il faut mesurer ce que le carton vient de faire. La liste Gal 1 à Gal 6 de la
page 157 exigeait la décomposition en composantes connexes comme un axiome ;
la liste G0 à G5 de la page 209 la déduit ; et la condition (v) ci-dessus
remplace toute la machinerie par une seule hypothèse d'existence, celle du
foncteur fibre. Trois états de la même axiomatique, dans un seul carton, et
c'est le dernier qui est passé dans SGA 1.

\subsection*{La proposition finale}

\textbf{Proposition.} Soit $u : C \to C'$ un foncteur \emph{exact}, $C$ et
$C'$ multigaloisiennes. Alors :

\begin{enumerate}
  \item[(i)] $u$ est \emph{fidèle} $\Longleftrightarrow$ $u$ est
    \emph{conservatif} $\Longleftrightarrow$ pour tout objet $X$ de $C$,
    $u(X) = \varnothing$ entraîne $X = \varnothing$.
  \item[(ii)] $u$ est \emph{pleinement fidèle} $\Longleftrightarrow$ pour tout
    objet $X$ de $C$, $u$ induit une \emph{bijection} de l'ensemble des
    sous-objets de $X$ sur l'ensemble des sous-objets de $u(X)$.
  \item[(iii)] $u$ est une \emph{équivalence} $\Longleftrightarrow$ $u$ est
    pleinement fidèle et tout objet de $C'$ est \emph{subordonné} à un objet de
    l'image.
\end{enumerate}

C'est sur cette proposition, dont la démonstration ne suit pas, que la cote
s'arrête ; la moitié inférieure de la page est nue.

\begin{quote}
\textbf{Ce que dit cette dernière page, et pourquoi elle est à sa place.} Ce
sont les critères sous lesquels les axiomes du domaine de ramification ont été
choisis, page 190, cent trente pages plus tôt dans le même carton. L'axiome
\textbf{R2} --- « l'application $X' \mapsto r(X')$ est une bijection de
l'ensemble des sous-objets directs de $X$ sur l'ensemble des parties ouvertes
et fermées de $r(X)$ » --- est le critère (ii). La proposition 2 de la page
192 --- « $X = \varnothing$ si et seulement si $r(X) = \varnothing$ » --- est
le critère (i). Et le critère (iii), la subordination, est le mot même sur
lequel toute la théorie des familles génératrices est bâtie aux pages 201 à
206.

La réalisation géométrique $r$ n'est pas un foncteur entre catégories
multigaloisiennes --- son but est la catégorie des schémas --- de sorte que la
proposition ne s'y applique pas littéralement. Mais c'est bien elle qui donne
aux axiomes R1 et R2 leur raison d'être : ils demandent de $r$ exactement les
propriétés qui, entre catégories multigaloisiennes, caractérisent la fidélité
et la pleine fidélité. Le carton s'achève donc sur les raisons de son propre
point de départ.
\end{quote}

\end{document}
